\chapter[1933 --- Morgan]{1933: Thomas Hunt Morgan}
\label{ch:1933_thomashuntmorgan}

\section*{Main Contribution}

\textit{Demonstrated that chromosomes carry the genes responsible for heredity, proving the chromosomal theory of inheritance.}

\section{Official Citation}

for his discoveries concerning the role played by the chromosomes in heredity

\section{Background and Historical Context}

Thomas Hunt Morgan was born on September 25, 1866, in Lexington, Kentucky. He studied science at the University of Kentucky and later at Johns Hopkins University, where he earned his doctorate in 1890 under the supervision of William Keith Brooks. Morgan's doctoral research was on the embryology of sea spiders (pycnogonids), a subject that reflected his early interest in marine biology and developmental anatomy. After a brief period of research in Europe, where he studied marine biology at the Naples Zoological Station---then one of the world's leading centers for marine research---Morgan joined the faculty of Bryn Mawr College in Pennsylvania, where he taught biology and conducted research on experimental embryology and regeneration.

In 1904, Morgan moved to Columbia University in New York, where he would conduct his most famous work on the genetics of the fruit fly \emph{Drosophila melanogaster}. At Columbia, Morgan built one of the most productive research laboratories in the history of biology, training a group of talented students---including Alfred Sturtevant, Calvin Bridges, and Hermann Muller---who would collectively establish the field of experimental genetics. In 1928, Morgan moved to the California Institute of Technology (Caltech), where he established the Division of Biology and built one of an influential biology departments in the world.

Morgan's Nobel Prize was awarded for his discoveries concerning the role played by the chromosome in heredity---a body of work that established the chromosomal theory of inheritance and transformed biology from a largely descriptive science into an experimental, molecular discipline. To understand the significance of Morgan's achievement, it is necessary to appreciate the state of genetic knowledge at the time.

The principles of heredity had been rediscovered in 1900, when three botanists---Hugo de Vries, Carl Correns, and Erich von Tschermak---independently rediscovered the work of Gregor Mendel, who had published his laws of inheritance in 1866. Mendel's laws---the law of segregation (each individual has two factors for each trait, and these factors separate during gamete formation) and the law of independent assortment (factors for different traits are inherited independently of each other)---provided a mathematical framework for understanding heredity, but they did not specify the physical nature of the ``factors'' (now called genes) or their location within the cell.

The hypothesis that genes were carried on chromosomes had been proposed by Walter Sutton and Theodor Boveri in 1902--1903, based on the observation that chromosomes occur in pairs (homologous chromosomes) that separate during meiosis, just as Mendel's factors separate during gamete formation. However, the Sutton-Boveri hypothesis remained unproven: no one had demonstrated a direct correspondence between a specific gene and a specific chromosome. It was Morgan who provided this proof, using the fruit fly \emph{Drosophila melanogaster} as his experimental organism.

\section{The Discovery and Contribution}

Morgan's approach to genetics was experimental and empirical. Rather than speculating about the nature of heredity, he set out to observe the patterns of inheritance in living organisms and to test whether these patterns were consistent with the chromosomal theory. His choice of \emph{Drosophila melanogaster} as an experimental organism was inspired: the fruit fly has a short generation time (about two weeks at room temperature), produces large numbers of offspring (a single female can lay hundreds of eggs), has only four pairs of chromosomes (which could be readily visualized under the microscope), and is easy to culture in the laboratory on simple media. These properties made it ideal for genetic analysis, and \emph{Drosophila} has remained one of the major model organisms in genetics for more than a century.

Morgan's breakthrough came in 1910, when he discovered a white-eyed male fly among a population of normally red-eyed flies. This white-eyed mutant was the first sex-linked mutation to be described, and it provided the first direct evidence that a gene is located on a specific chromosome. When Morgan crossed the white-eyed male with a red-eyed female, all the offspring had red eyes (demonstrating that red eyes were dominant over white eyes). When the F1 offspring were crossed with each other, the F2 generation showed the expected 3:1 ratio of red-eyed to white-eyed flies, but---crucially---all the white-eyed flies were male. This pattern of inheritance was consistent with the hypothesis that the white-eye gene was located on the X chromosome (one of the two sex chromosomes), because males have only one X chromosome (inherited from their mother) and therefore express any gene located on the X chromosome, whether dominant or recessive.

This discovery---that a specific gene was located on a specific chromosome---was the first experimental proof of the Sutton-Boveri chromosomal theory of inheritance, and it opened up an entirely new era in genetics. Building on this initial discovery, Morgan and his students conducted a systematic analysis of the inheritance patterns of hundreds of mutations in \emph{Drosophila}, including mutations affecting eye color, wing shape, bristle pattern, body color, and many other characters.

They showed that genes located on the same chromosome tend to be inherited together (a phenomenon called ``linkage''), but that linkage is not absolute: genes on the same chromosome can be separated by crossing over (the exchange of genetic material between homologous chromosomes during meiosis). The frequency of crossing over between two genes is proportional to the distance between them---a principle that Sturtevant used to construct the first genetic map in 1913. Sturtevant, who was then a 19-year-old undergraduate working in Morgan's laboratory, mapped the relative positions of six genes on the X chromosome of \emph{Drosophila}, demonstrating that genes are arranged in a linear order on chromosomes.

The genetic map was a conceptual and practical breakthrough. It demonstrated that genes are arranged in a linear order on chromosomes, like beads on a string, and it provided a tool for predicting the patterns of inheritance of multiple genes simultaneously. The construction of genetic maps became the foundation of classical genetics, and it remains an important tool in modern genomics, where it has been extended from the analysis of individual genes to the mapping of entire genomes.

Morgan also contributed to the understanding of the mechanism of crossing over and the nature of the gene. He showed that crossing over involved the physical exchange of segments between homologous chromosomes, and he provided evidence that genes were discrete units that could be separated from each other by recombination. These observations contributed to the development of the ``beads on a string'' model of gene arrangement on chromosomes, which prevailed for decades before being superseded by the molecular understanding of genes as segments of DNA.

Morgan's work on \emph{Drosophila} genetics was summarized in his monumental book, \emph{The Mechanism of Mendelian Heredity} (1915, co-authored with Muller, Sturtevant, and Bridges), which became the standard reference in the field and established \emph{Drosophila} as the premier model organism for genetic research. The book was notable for its comprehensive coverage of genetic principles, its clear presentation of experimental evidence, and its integration of genetic theory with cytological observations of chromosomes.

Morgan received the Nobel Prize in Physiology or Medicine in 1933 ``for his discoveries concerning the role played by the chromosome in heredity.'' The award recognized not only Morgan's specific discoveries but also the general paradigm of experimental genetics that he established. Morgan was the only member of the ``Drosophila school'' to receive the Nobel Prize, but the influence of his students and collaborators extended far beyond his own laboratory.

\section{Impact on Modern Medicine}

Morgan's discoveries have had a transformative impact on modern medicine, from the diagnosis of genetic diseases to the development of gene therapy and personalized medicine. The chromosomal theory of inheritance that Morgan established is the foundation of medical genetics, and the tools of genetic analysis that he pioneered---linkage mapping, genetic screening, and mutation detection---are now standard components of clinical practice.

The identification of disease-causing genes through linkage analysis and positional cloning is a direct application of the mapping techniques that Morgan and Sturtevant developed. The discovery of the genes responsible for cystic fibrosis (1989), Huntington's disease (1993), sickle cell anemia (1978), and thousands of other genetic disorders has been made possible by the principles of genetic linkage and chromosomal mapping that Morgan elucidated. The Human Genome Project, which sequenced the entire human genome and was completed in 2003, built upon the genetic mapping techniques that originated in Morgan's \emph{Drosophila} laboratory. The reference human genome, which serves as the template for all subsequent genomic analyses, is a direct intellectual descendant of the genetic maps that Morgan's students constructed a century ago.

In modern oncology, the chromosomal theory of inheritance has been extended to the understanding of cancer as a disease of the genome. The discovery that cancer is caused by mutations in specific genes---oncogenes (which promote cell growth when activated) and tumor suppressor genes (which restrain cell growth when active and cause cancer when inactivated)---that are located on specific chromosomes is a direct application of Morgan's chromosomal theory. The identification of cancer-predisposing mutations (such as BRCA1 and BRCA2 for breast cancer, and APC for colorectal cancer) has enabled the development of genetic screening programs and targeted therapies that are transforming the treatment of cancer. The recognition that cancer is fundamentally a genetic disease---caused by the accumulation of mutations in critical genes---is one of the major conceptual advances in modern oncology, and it has its roots in Morgan's demonstration that genes are carried on chromosomes.

Gene therapy, which aims to treat or cure genetic diseases by introducing functional copies of defective genes, is a direct descendant of Morgan's work on the relationship between genes and chromosomes. The development of CRISPR-Cas9 gene editing technology, which allows precise modification of genes at specific chromosomal locations, is the most recent and most powerful extension of the genetic principles that Morgan established. The potential to correct genetic defects at their source---to rewrite the code of heredity---represents the ultimate fulfillment of the vision that Morgan's work initiated. Clinical trials of CRISPR-based therapies for sickle cell anemia, beta-thalassemia, and certain forms of inherited blindness are already underway, offering hope that Morgan's legacy will soon translate into cures for some of the most devastating genetic diseases.

Morgan's choice of \emph{Drosophila} as a model organism also established the paradigm of model organism research that dominates modern biology. The use of genetically tractable organisms---yeast (\emph{Saccharomyces cerevisiae}), nematode worms (\emph{Caenorhabditis elegans}), fruit flies (\emph{Drosophila melanogaster}), and mice (\emph{Mus musculus})---to study fundamental biological processes is now a cornerstone of biomedical research, and the insight that genetic analysis can be most effectively conducted in organisms with short generation times, large offspring numbers, and tractable genomes is Morgan's enduring legacy.

\section{Legacy}

Thomas Hunt Morgan's legacy is that of a scientist who transformed biology from a largely descriptive science into an experimental, quantitative discipline. By demonstrating that genes are carried on chromosomes and that they can be mapped and analyzed with mathematical precision, Morgan laid the foundation for the modern science of genetics.

Morgan was a gifted experimentalist and a charismatic teacher. The ``fly room'' at Columbia University, where Morgan and his students conducted their legendary experiments, was one of the most productive laboratories in the history of biology. The ``Drosophila school'' that Morgan established produced several Nobel laureates (including Muller and, indirectly, Edward B. Lewis) and trained a generation of geneticists who spread Morgan's methods and ideas throughout the world. The fly room was also notable for its informal and collaborative atmosphere, which encouraged creative thinking and the free exchange of ideas---a model for research laboratories that persists to this day.

After receiving the Nobel Prize, Morgan continued to contribute to biology at Caltech, where he established the first formal biology department at a research university and fostered the interdisciplinary approach to biology that is now standard. His vision of biology as a unified science, integrating genetics, biochemistry, cell biology, and developmental biology, anticipated the modern systems biology approach that seeks to understand biological processes at all levels of organization. Morgan also served as the president of the National Academy of Sciences from 1927 to 1931, and he was a tireless advocate for the support of basic research by government and private funding agencies.

Morgan was also a thoughtful commentator on the relationship between science and society. He was concerned about the potential misuse of genetic knowledge (particularly in the context of the eugenics movement, which was gaining strength in the early twentieth century), and he argued that genetic information should be used to help individuals rather to make judgments about the relative worth of different groups of people. His ethical stance anticipated the debates about genetic privacy, genetic discrimination, and the responsible use of genetic information that continue to shape public policy in the era of genomic medicine. The Genetic Information Nondiscrimination Act (GINA) of 2008, which prohibits discrimination based on genetic information in employment and health insurance, is a modern reflection of the ethical principles that Morgan championed.

The chromosomal theory of inheritance that Morgan established is now more than a century old, but its influence continues to grow. Every time a geneticist maps a disease gene to a specific chromosomal location, every time a physician orders a genetic test to diagnose a hereditary condition, and every time a researcher uses CRISPR to edit a gene, they are building upon the foundations that Thomas Hunt Morgan laid in his fly room at Columbia University more than a hundred years ago. Morgan died in Pasadena, California, on December 4, 1945, but his scientific legacy endures in every advance in genetics and genomics that has occurred since his pioneering work.


\section{Modern Developments}

Morgan's chromosomal theory has evolved into modern genomics. The Human Genome Project (2003) sequenced all 3 billion base pairs of human DNA. CRISPR-Cas9 gene editing (2020 Nobel Prize) enables precise modification of genes on chromosomes. Chromosome painting and FISH (fluorescence in situ hybridization) visualize individual chromosomes in clinical diagnostics. Single-cell sequencing reveals chromosomal abnormalities in individual cells, transforming cancer diagnosis.
